\chapter{La première nuit}
\sessioninfo{Session 9}{Mardi 27 octobre 2026}

Le calme revint peu à peu sur la rivière. Une araignée géante gisait sur la rive, tandis que le courant avait déjà emporté la seconde. Kael retrouva l'une de ses flèches, intacte, plantée dans un arbre ; l'autre avait disparu avec la créature tombée à l'eau. Non loin de là, Obélik récupéra sa hachette dans la végétation.

La toile barrait toujours le passage au-dessus de l'eau. Morrÿs tendit une main et projeta un \textit{Fire Bolt} dans les filaments. Le feu courut presque aussitôt sur la soie, qui se recroquevilla avant de disparaître en quelques instants. La voie était de nouveau libre.

Le groupe profita de cette courte halte pour s'assurer que personne n'avait été gravement blessé. Obélik portait encore les traces de la morsure reçue pendant l'affrontement, mais le goliath affirma avoir connu bien pire. Azaka, elle aussi, avait été mordue. Pourtant, sa plaie ne saignait déjà plus et semblait s'être refermée d'elle-même. Cette étrange guérison n'échappa pas à ses compagnons. La guide leur accorda seulement une réponse évasive : sa lignée portait de vieilles cicatrices et du vieux sang, et les dangers de Chult ne lui étaient pas étrangers.

Kael s'intéressa plutôt au venin des araignées. Azaka confirma qu'il était efficace, mais expliqua que ce qu'ils pourraient prélever ne serait encore qu'une essence brute, inutilisable sans préparation. Morrÿs rappela que la possession et le commerce de poisons pouvaient attirer de sérieux ennuis sur la Côte des Épées, même si l'usage personnel semblait ouvrir la porte à quelques interprétations commodes.

Kael examina les mandibules de la créature restée sur la rive et repéra les glandes qu'il cherchait. Il tenta de les dégager avec précaution, mais comprit rapidement qu'il risquait surtout de les endommager. Plutôt que de gâcher ce qui restait, il abandonna l'opération. Le groupe repartit donc sans venin.

\scenebreak

Ils remirent le canoë à l'eau et reprirent leur descente de la rivière. À mesure que la lumière déclinait, l'air devint plus lourd et plus humide. Azaka leva les yeux vers le ciel, puis leur conseilla de gagner la rive sans tarder : la pluie arrivait.

Le canoë fut tiré hors de l'eau et retourné afin qu'il ne puisse ni dériver ni se remplir. Le groupe installa le récupérateur d'eau, rassembla du bois et parvint à faire naître un feu. Ils venaient à peine de terminer leurs préparatifs lorsque le ciel s'ouvrit.

Pendant plusieurs minutes, la pluie tropicale s'abattit sur le campement avec une violence stupéfiante. L'eau ruissela sur les feuillages, creusa des flaques dans la terre et étouffa le feu en quelques instants. Le récupérateur, en revanche, se remplit rapidement et leur procura assez d'eau pour remplir toutes les gourdes.

Azaka avait pris soin de conserver du bois sec. Elle gratta la couche détrempée du sol jusqu'à trouver de quoi protéger une nouvelle flamme, puis ralluma le foyer. Les branches humides furent disposées tout autour afin qu'elles sèchent peu à peu. Destos purifia l'eau recueillie pendant que chacun achevait d'installer le camp.

Obélik partit couper quelques branches supplémentaires. Le goliath observa avec curiosité les arbres de la jungle, si différents de ceux qu'il avait connus : des troncs noueux se mêlaient aux lianes, des racines plongeaient profondément à la recherche d'eau et les feuilles se disputaient la lumière bien au-dessus de leurs têtes. Par endroits, il devenait difficile de distinguer où finissait un arbre et où commençait le suivant.

Ils partagèrent leur repas autour du feu. C'était leur premier véritable campement dans la jungle, loin des murs de Port Nyanzaru et de la sécurité relative de l'\textit{Effleureuse}. Avant de s'allonger, Kael traça autour d'eux un périmètre magique. Si une présence étrangère le franchissait, une cloche devait donner l'alerte sur une note soigneusement choisie.

\begin{narrateur}
Un la, réglé à quatre cent quarante hertz. Même au cœur de la jungle, certaines choses ne souffrent aucune approximation.
\end{narrateur}

\scenebreak

Obélik prit le premier tour de garde. La jungle bruissait tout autour de lui, pleine de froissements, de cris lointains et de mouvements impossibles à identifier, mais rien ne vint troubler le camp. Au terme de son quart, il réveilla Destos et alla s'allonger près du feu.

La cloche magique retentit pendant la garde du clerc. Destos chercha aussitôt la trace d'une présence surnaturelle, sans rien percevoir de tel. En observant plus attentivement les abords du camp, il finit par découvrir le responsable : un petit ragondin s'approchait prudemment des restes du repas.

Destos tenta de transformer l'intrus en provisions pour les jours suivants. Une flamme sacrée s'abattit près de l'animal, qui détala dans la végétation avant qu'il puisse recommencer. Le clerc dut se contenter de réveiller Morrÿs, sans viande supplémentaire à ajouter aux réserves.

Le quart du nain se déroula sans incident. Il scruta les alentours et surveilla même la terre autour du campement, mais ne remarqua rien d'inquiétant. Lorsque son tour s'acheva, il réveilla Kael.

Le rôdeur sentit bientôt l'humidité augmenter de nouveau. Une seconde averse approchait, plus durable que la première. Il mit à l'abri ce qui pouvait encore l'être, protégea le bois sec sous le canoë retourné et se prépara au déluge.

La pluie tomba longtemps, gonflant les flaques et faisant courir de petits ruisseaux de boue vers la rivière. Malgré l'inconfort, aucun danger ne se présenta. Lorsque l'averse finit par faiblir, Kael réveilla Azaka pour les deux dernières heures de garde.

La guide constata qu'il n'avait pas profité de toute cette eau gratuite pour se laver. Kael lui répondit qu'il avait eu d'autres préoccupations. Elle jugea cela regrettable, puis prit sa place tandis qu'il rejoignait les autres.

Une fois encore, la cloche magique finit par se manifester, au grand agacement d'Azaka. Quelque petite créature avait dû franchir le périmètre avant de disparaître dans la jungle. Rien d'autre ne troubla la fin de la nuit.

Au matin, le groupe s'éveilla au milieu d'un campement trempé. La pluie avait cessé, la rivière poursuivait sa course et le canoë les attendait sur la rive. Leur première nuit dans la jungle était derrière eux.
